\documentclass[UTF8,a4paper,11pt,openany]{ctexbook}
\usepackage{../common/statlearnbook}
\SLSetBookMetadata{第 1 册}{统计学习方法讲义 v2.6.0}{数学基础与统计学习基本理论}
\begin{document}
\setcounter{page}{4}
\setcounter{chapter}{1}
\setcounter{figure}{0}
\pagestyle{headings}
\chaptermark{数学语言、符号约定与学习路线}
\begin{geometrybox}
\textbf{几何解释。} 向量$\boldsymbol{x}\in\R^d$可视为$d$维空间中的点，也可视为从原点指向该点的箭头；矩阵$X$把$n$个样本按行排列。维数核对相当于检查几何变换的输入空间与输出空间能否衔接。
\end{geometrybox}
\begin{probabilitybox}
事件是样本空间的子集，事件指示函数是取值为$0$或$1$的随机变量。其平均值记录事件在样本中的发生比例，并为后续“期望等于概率”的关系准备统一语言。
\end{probabilitybox}
\begin{optimizationbox}
\textbf{优化解释。} 参数$\theta$是搜索对象，允许取值的集合是可行域，评价参数优劣的函数是目标函数。先声明变量与条件，才能判断一次更新是否仍在可行域内。
\end{optimizationbox}
% v2.6.0 figure UID: FIG-P020-01
% Image-informed reverse drawing: preserve the dependency graph while making
% the main chain, audit return, text hierarchy, and endpoint clearance explicit.
\tikzset{slfig-FIG-P020-01/.style={font=\fontsize{10.0pt}{12.0pt}\selectfont,every path/.append style={line cap=round,line join=round}}}
\pgfkeysifdefined{/tikz/alt}{}{\tikzset{alt/.code={}}}
\begin{figure}[H]\centering\label{schema:V1-C01:CH-17}
\begin{tikzpicture}[slfig-FIG-P020-01,
  node distance=6.2mm,
  stage/.style={rounded corners=1.7mm,draw=SLBlue,line width=.90pt,
    text width=29mm,minimum height=13.2mm,align=center,inner xsep=3.2pt,inner ysep=3pt,
    font=\fontsize{10.5pt}{12.6pt}\selectfont,text=SLInk},
  alt={对象—关系—结论：数学语言从对象声明到任务陈述的依赖关系。每一条箭头都表示右侧内容使用左侧定义。}
]
  \node[stage,fill=SLBlueBG] (object) {{\bfseries 对象声明}\\[-.6pt]{\fontsize{10.0pt}{12.0pt}\selectfont 集合、类型与维数}};
  \node[stage,fill=SLTealBG,right=of object] (relation) {{\bfseries 关系与映射}\\[-.6pt]{\fontsize{10.0pt}{12.0pt}\selectfont 定义域%
    \tikz[baseline=-.55ex]{
      \path[use as bounding box] (0,-.75mm) rectangle (7mm,.75mm);
      \draw[-{Stealth[length=1.55mm,width=1.05mm]},draw=SLInk,line width=.72pt]
        (1.05mm,0)--(5.95mm,0);
    }%
    值域}};
  \node[stage,fill=SLBlueBG,right=of relation] (logic) {{\bfseries 运算与逻辑}\\[-.6pt]{\fontsize{10.0pt}{12.0pt}\selectfont 复合、量词与约束}};
  \node[stage,fill=SLTealBG,right=of logic] (task) {{\bfseries 可核验任务}\\[-.6pt]{\fontsize{10.0pt}{12.0pt}\selectfont 输入、输出与判据}};
  \begin{scope}[on background layer]
    \node[fit=(object)(task),rounded corners=2.4mm,fill=SLBlue!3,draw=none,
      inner xsep=2.6mm,inner ysep=2.8mm] {};
  \end{scope}
  \foreach \a/\b in {object/relation,relation/logic,logic/task}
    \draw[-{Stealth[length=2.25mm,width=1.45mm]},draw=SLBlue,line width=.98pt,
      shorten <=1.5mm,shorten >=1.5mm]
      (\a.east)--(\b.west);
  \draw[-{Stealth[length=2.0mm,width=1.25mm]},draw=SLGray,line width=.58pt,
      dash pattern=on 2.4pt off 2.0pt,shorten <=1.5mm,shorten >=1.5mm]
    (task.south)--++(0,-9mm)-| node[pos=.50,below=2.0pt,fill=white,inner sep=1pt,
      font=\fontsize{10.0pt}{12.0pt}\selectfont,text=SLGray] {逆向核对：任务所用定义逐项返回检查}
    (object.south);
\end{tikzpicture}
\caption{数学语言从对象声明到任务陈述的依赖关系。每一条箭头都表示右侧内容使用左侧定义。}\label{fig:V1-C01-language-flow}
\end{figure}

使用\cref{fig:V1-C01-language-flow}时应从任务端逆向核对：对象、取值域、映射和目标只要有一项未定义，就先补全声明再计算；箭头记录使用关系，不是可逆蕴含。
\begin{example}[三样本数据的类型核对]
设$X=\begin{bmatrix}1&0\\2&1\\0&1\end{bmatrix}\in\R^{3\times2}$，$\boldsymbol{w}=(2,-1)^{\mathsf T}\in\R^2$，标签$\boldsymbol{y}=(1,1,-1)^{\mathsf T}$。求$X\boldsymbol{w}$，并核对结果维数。
\end{example}
\end{document}
